\documentclass[12pt,a4paper]{article}

% === PACKAGES ===
\usepackage[utf8]{inputenc}
\usepackage[T1]{fontenc}
\usepackage{amsmath,amssymb,amsfonts}
\usepackage{graphicx}
\usepackage{geometry}
\geometry{margin=2.5cm}
\usepackage{hyperref}
\usepackage{xcolor}
\usepackage{booktabs}
\usepackage{float}
\usepackage{listings}
\usepackage{paralist}
\usepackage{caption}
\usepackage{subcaption}
\usepackage{titlesec}
\usepackage{fancyhdr}
\usepackage{colortbl}
\usepackage{url}

% === COLORS ===
\definecolor{accent}{RGB}{123,47,247}       % Electric violet
\definecolor{accentCyan}{RGB}{0,217,255}   % Neon cyan
\definecolor{accentPink}{RGB}{255,45,120}  % Neon pink
\definecolor{darkBg}{RGB}{13,11,30}         % Deep midnight
\definecolor{cardBg}{RGB}{26,23,48}        % Card surface
\definecolor{cardBorder}{RGB}{45,43,80}    % Subtle border
\definecolor{muted}{RGB}{138,138,168}      % Muted text
\definecolor{success}{RGB}{0,229,160}       % Neon green
\definecolor{codebg}{RGB}{22,20,40}        % Code background
\definecolor{heading}{RGB}{230,230,255}    % Heading color

% === TITLE PAGE ===
\graphicspath{{./}}
\hypersetup{
    colorlinks=true,
    linkcolor=accentCyan,
    urlcolor=accentCyan,
    pdfborder={0 0 0}
}

% === CUSTOM COMMANDS ===
\newcommand{\pingme}{PingMe\xspace}
\newcommand{\code}[1]{\texttt{#1}}
\newcommand{\file}[1]{\texttt{\detokenize{#1}}}
\newcommand{\pkg}[1]{\texttt{#1}}
\newcommand{\figref}[1]{Figure~\ref{#1}}
\newcommand{\tabref}[1]{Table~\ref{#1}}

% === LISTINGS STYLE ===
\lstdefinestyle{latex}{
    language=Tex,
    backgroundcolor=\color{codebg},
    basicstyle=\ttfamily\small\color{accentCyan},
    commentstyle=\color{muted}\itshape,
    keywordstyle=\color{accentPink}\bfseries,
    stringstyle=\color{success},
    numbers=none,
    frame=single,
    framerule=0.5pt,
    rulecolor=\color{cardBorder},
    breaklines=true,
    breakatwhitespace=true,
    xleftmargin=0.5cm,
    xrightmargin=0.5cm,
    aboveskip=8pt,
    belowskip=8pt,
}
\lstdefinestyle{js}{
    language=JavaScript,
    backgroundcolor=\color{codebg},
    basicstyle=\ttfamily\small\color{accentCyan},
    commentstyle=\color{muted}\itshape,
    keywordstyle=\color{accentPink}\bfseries,
    stringstyle=\color{success},
    numbers=left,
    numbersep=8pt,
    frame=single,
    framerule=0.5pt,
    rulecolor=\color{cardBorder},
    breaklines=true,
    breakatwhitespace=true,
    xleftmargin=0.5cm,
    xrightmargin=0.5cm,
    aboveskip=8pt,
    belowskip=8pt,
}

% === SECTION FORMATTING ===
\titleformat{\section}{\Large\bfseries\color{heading}}{}{0pt}{}
\titleformat{\subsection}{\large\bfseries\color{heading}}{}{0pt}{}
\titleformat{\subsubsection}{\normalsize\bfseries\color{heading}}{}{0pt}{}
\titlespacing*{\section}{0pt}{1.2cm}{0.4cm}
\titlespacing*{\subsection}{0pt}{0.8cm}{0.3cm}
\titlespacing*{\subsubsection}{0pt}{0.5cm}{0.2cm}

% === HEADERS / FOOTERS ===
\pagestyle{fancy}
\fancyhf{}
\fancyhead[L]{\color{muted}\small\pingme{} Project Documentation}
\fancyhead[R]{\color{muted}\small\thepage}
\fancyfoot[C]{\color{muted}\small Draft --- April 2026}
\def\headrule{\color{cardBorder}}
\def\footrule{\color{cardBorder}}

% === DOCUMENT ===
\begin{document}

% --- COVER PAGE ---
\thispagestyle{empty}
\vspace*{2cm}
\begin{center}
    \rule{0.95\linewidth}{0.5pt}
    \vspace{1cm}

    {\Huge\bfseries\color{accentCyan}\pingme{}}
    \vspace{0.4cm}

    {\LARGE\bfseries\color{heading}Project Report}
    \vspace{0.8cm}

    {\Large\sffamily\color{offWhite}Real-Time Messaging Application}
    \vspace{0.6cm}

    {\large\sffamily\color{muted}A Full-Stack Implementation with
    Socket.io, WebRTC, MongoDB, and React}

    \vspace{2cm}
    \rule{0.95\linewidth}{0.5pt}
    \vspace{2cm}

    \begin{tabular}{rl}
        \textbf{Status:} & \color{success}\textbf{In Progress} \\
        \textbf{Version:} & 1.0.0 \\
        \textbf{Date:}   & April 2026 \\
        \textbf{Branch:} & dev \\
    \end{tabular}
\end{center}
\newpage

% --- TABLE OF CONTENTS ---
\fancyhead[L]{\color{muted}\small Contents}
\tableofcontents
\newpage

% --- SECTION 1: INTRODUCTION ---
\section{Introduction}
\label{sec:introduction}

\pingme{} is a full-stack real-time messaging application with a sci-fi neural network
user interface (dark backgrounds, neon gradient accents, cyberpunk theme). The application
provides one-on-one chat, friend management, file sharing, and voice/video calling.

The project follows a client--server architecture with two independent Node.js packages:
a React frontend and an Express backend, communicating via REST API and Socket.io for
real-time events.

\subsection{Project Goals}
\begin{compactitem}
    \item Deliver a real-time messaging experience comparable to modern chat applications.
    \item Support rich media sharing including images, documents, and file attachments.
    \item Provide voice and video calling via WebRTC peer-to-peer connections.
    \item Maintain security through JWT authentication with httpOnly cookies.
    \item Offer an immersive cyberpunk UI with neon aesthetics and smooth animations.
\end{compactitem}

\subsection{Technology Stack}
\label{sec:tech-stack}

\begin{table}[H]
    \centering
    \caption{Technology Stack Overview}
    \begin{tabular}{llp{6cm}}
        \toprule
        \textbf{Layer}    & \textbf{Tech}           & \textbf{Role} \\
        \midrule
        Frontend         & React 19 + Vite 7       & SPA, state management \\
        Styling         & Tailwind CSS v4          & CSS-first theming \\
        Real-time       & Socket.io-client         & Bidirectional events \\
        HTTP Client     & Axios                   & REST calls + 401 refresh \\
        Backend         & Express 5 (Node.js ESM) & REST API, Socket.io server \\
        Database        & MongoDB + Mongoose 8    & Document store \\
        Auth           & JWT + bcrypt            & Access + refresh tokens \\
        File Upload     & Multer                  & Multipart handling \\
        \bottomrule
    \end{tabular}
\end{table}

% --- SECTION 2: SYSTEM ARCHITECTURE ---
\section{System Architecture}
\label{sec:architecture}

\subsection{Overall Architecture}

\pingme{} uses a client--server model with a three-layer architecture as shown in
\figref{fig:architecture}.

\begin{figure}[H]
    \centering
    \fbox{
        \begin{minipage}{0.9\linewidth}
            \centering
            \vspace{0.5cm}
            \textbf{Client (React 19 + Vite)} \\
            \small AuthContext, SocketContext, ChatArea, MessageList, MessageBubble, MessageInput \\
            \vspace{0.4cm}
            \hrule
            \vspace{0.4cm}
            \textbf{Server (Express 5 + Socket.io 4)} \\
            \small REST Controllers, Socket.io Handler, Multer Upload, JWT Middleware \\
            \vspace{0.4cm}
            \hrule
            \vspace{0.4cm}
            \textbf{Database (MongoDB)} \\
            \small Users, Messages \\
            \vspace{0.5cm}
        \end{minipage}
    }
    \caption{Three-Layer Architecture}
    \label{fig:architecture}
\end{figure}

The frontend communicates with the backend via two channels:
\begin{compactitem}
    \item \textbf{REST API} (Axios): Authentication, message retrieval, file upload.
    \item \textbf{Socket.io}: Real-time message delivery, typing indicators,
          friend status, call signaling, reactions.
\end{compactitem}

\subsection{Project Structure}
\label{sec:project-structure}

\begin{figure}[H]
    \centering
    \fbox{
        \begin{minipage}{0.92\linewidth}
            \small
            \raggedright
            PingMe/ \\
            \hspace{1em}\file{client/} \\
            \hspace{2em}\file{src/} \\
            \hspace{3em}\file{components/} \\
            \hspace{4em}\file{layout/} --- AuthLayout, ChatArea, Header, Sidebar, TopNavBar \\
            \hspace{4em}\file{chat/}     --- MessageBubble, MessageInput, MessageList, EmojiPicker \\
            \hspace{4em}\file{ui/}        --- Avatar, Button, Input, SearchBar, ValidationMessage \\
            \hspace{3em}\file{context/}  --- AuthContext, SocketContext \\
            \hspace{3em}\file{pages/}    --- Login, Register, Chat \\
            \hspace{3em}\file{config/}    --- api.js (Axios instance) \\
            \hspace{3em}\file{socket.js}   --- Socket.io singleton \\
            \hspace{2em}\file{server/} \\
            \hspace{3em}\file{index.js}   --- Express app + HTTP server + Socket.io \\
            \hspace{3em}\file{config/}    --- db.js (Mongoose connection) \\
            \hspace{3em}\file{models/}     --- User, Message \\
            \hspace{3em}\file{controllers/} --- auth, message, user \\
            \hspace{3em}\file{services/}  --- auth.service.js \\
            \hspace{3em}\file{middlewares/} --- auth.middleware, upload.middleware \\
            \hspace{3em}\file{routers/}   --- auth, message, user routes \\
            \hspace{3em}\file{socket/}    --- socketHandler.js \\
            \hspace{3em}\file{uploads/}   --- (gitignored) \\
        \end{minipage}
    }
    \caption{Project Directory Structure}
    \label{fig:structure}
\end{figure}

% --- SECTION 3: AUTHENTICATION ---
\section{Authentication System}
\label{sec:auth}

The authentication system uses a dual-token strategy with httpOnly cookies for security.

\subsection{Token Strategy}

\begin{compactitem}
    \item \textbf{Access Token} (15 minutes): Short-lived JWT stored in memory.
      Verified via \code{authMiddleware}.
    \item \textbf{Refresh Token} (7 days): Long-lived JWT in httpOnly cookie.
      Used silently when the access token expires.
\end{compactitem}

The Axios interceptor detects 401 responses, calls \code{POST /auth/refresh},
updates both tokens, and retries the failed request automatically.

\subsection{Auth Middleware Logic}
\label{sec:auth-middleware}

The \code{authMiddleware} checks tokens in this priority order:

\begin{lstlisting}[style=js]
// 1. Check httpOnly cookie first (refresh token flow)
const refreshToken = req.cookies?.refreshToken;
if (refreshToken) {
    const decoded = jwt.verify(refreshToken, ACCESS_TOKEN_SECRET);
    req.user = decoded;
    return next();
}
// 2. Fall back to Authorization: Bearer header
const authHeader = req.headers.authorization;
if (authHeader?.startsWith('Bearer ')) {
    const token = authHeader.slice(7);
    const decoded = jwt.verify(token, ACCESS_TOKEN_SECRET);
    req.user = decoded;
    return next();
}
return res.status(401).json({ error: 'Unauthorized' });
\end{lstlisting}

Note: the User model uses \code{select: false} for the \code{password} field.
Use \code{.select('+password')} when comparing passwords.

% --- SECTION 4: REAL-TIME MESSAGING ---
\section{Real-Time Messaging}
\label{sec:messaging}

\subsection{Socket.io Setup}
\label{sec:socket-setup}

The socket singleton (\file{client/src/socket.js}) is created with \code{autoConnect: false}.
\code{SocketProvider} auto-connects on mount and registers the user:

\begin{lstlisting}[style=js]
// Client emits on Chat.jsx mount:
socket.emit('register_user', userId);

// Server maintains in-memory map:
const onlineUsers = new Map(); // userId -> socketId
\end{lstlisting}

After registration the server broadcasts online friends' status to the client.

\subsection{Message Flow}
\label{sec:message-flow}

Messages follow this flow:

\begin{enumerate}
    \item User sends message in \code{MessageInput}.
    \item Client emits \code{send\_message} via Socket.io.
    \item Server saves message to MongoDB (collection: \code{messages}).
    \item Server relays message to recipient via \code{receive\_message}.
    \item Server confirms to sender via \code{message\_sent}.
    \item Both sides update their message list state.
\end{enumerate}

\subsection{Status Indicators}
\label{sec:status-indicators}

Typing state is managed via two Socket.io events:
\code{typing} / \code{stop\_typing} (client-to-server) and
\code{user\_typing} / \code{user\_stopped\_typing} (server-to-client).

Read receipts use \code{mark\_messages\_read} (client-to-server) and
\code{messages\_were\_read} (server-to-client).

% --- SECTION 5: FRIEND MANAGEMENT ---
\section{Friend Management}
\label{sec:friends}

Users can search, send requests, accept requests, and browse a friend list.
All friend operations are exposed via REST API endpoints:

\begin{table}[H]
    \centering
    \caption{Friend Management API Endpoints}
    \begin{tabular}{llp{7cm}}
        \toprule
        \textbf{Method} & \textbf{Endpoint}             & \textbf{Description} \\
        \midrule
        GET    & \code{/api/users}              & List all users except self \\
        GET    & \code{/api/users/friends}     & Get friend list \\
        GET    & \code{/api/users/requests}    & Get pending friend requests \\
        GET    & \code{/api/users/search?q=}   & Search users by name \\
        POST   & \code{/api/users/request}      & Send a friend request \\
        POST   & \code{/api/users/accept}       & Accept a friend request \\
        \bottomrule
    \end{tabular}
\end{table}

The User model stores \code{friends: []} and \code{friendRequests: []}
as arrays of user IDs. A "reject friend request" button exists in the UI
but the backend route has not been implemented yet.

% --- SECTION 6: MEDIA AND FILE SHARING ---
\section{Media and File Sharing}
\label{sec:media}

The media sharing module allows users to send images and file attachments,
preview media, react to messages, and browse a conversation gallery.

\subsection{Image Upload}
\label{sec:image-upload}

Images are uploaded via \code{POST /api/messages/upload} with \code{FormData}.
The backend uses Multer with a 25\,MB file size limit and stores files in
\file{server/uploads/} (gitignored). File naming convention:
\code{{userId}\_{timestamp}\_{originalname}}.

Supported image types are validated in the Multer \code{fileFilter}.
Preview is generated client-side with \code{URL.createObjectURL()} before upload.

Acceptance criteria for image upload:
\begin{compactitem}
    \item Native file picker (image/*), max 10\,MB per image.
    \item Preview before sending, progress indicator during upload.
    \item Multi-image gallery: up to 5 images per message.
    \item Retry button on failure.
    \item Clipboard paste support.
\end{compactitem}

\subsection{File Attachments}
\label{sec:file-attachments}

Non-image files are displayed as file cards (icon + filename + size).
Supported formats: PDF, DOC, DOCX, XLS, XLSX, PPT, PPTX, ZIP, RAR, TXT, CSV.
Maximum file size: 25\,MB.

Clicking a file card triggers browser download via the static file URL:
\url{http://localhost:3001/uploads/{filename}}.

\subsection{Image Lightbox}
\label{sec:lightbox}

Tapping an image opens a fullscreen overlay with a dark backdrop at 90\% opacity.
Controls: close button (top-right), click backdrop to close,
swipe down to close (mobile), pinch-to-zoom support.

The Message model has been extended with:
\code{attachments: [\{ type, url, filename, size \}]}.

\subsection{Message Reactions}
\label{sec:reactions}

Six quick reactions are available: \texttt{:+1: :heart: :joy: :fire: :clap: :poop:}.
Triggered by long-press (500\,ms) or hover. Reactions emit via Socket.io:

\begin{lstlisting}[style=js]
Client --> Server: add_reaction    { messageId, emoji, userId }
Client --> Server: remove_reaction { messageId, emoji, userId }
Server --> Client: reaction_added  { messageId, emoji, userId }
Server --> Client: reaction_removed { messageId, emoji, userId }
\end{lstlisting}

The Message model now includes: \code{reactions: [\{ emoji, userId \}]}.

Tapping a reaction badge shows the list of users who reacted.
Reacting with the same emoji again toggles the reaction off.

\subsection{Media Gallery}
\label{sec:gallery}

Clicking the gallery icon in the chat header opens a grid panel showing
all images in the conversation (newest first). Clicking a thumbnail
opens the lightbox at that image. Empty state message:
``Chua co anh nao duoc chia se''.

% --- SECTION 7: VOICE AND VIDEO CALLS ---
\section{Voice and Video Calls}
\label{sec:calls}

Voice and video calling uses WebRTC for peer-to-peer connections and
Socket.io for signaling (offer/answer/ICE exchange).

\subsection{STUN Configuration}
\label{sec:stun}

The STUN server is configured using Google's public server:

\begin{lstlisting}[style=js]
const iceServers = [
    { urls: 'stun:stun.l.google.com:19302' }
    // TURN server required for production (symmetric NAT)
];
\end{lstlisting}

If no ICE candidates are found within 10 seconds,
a ``Khong the ket noi'' error is shown. A 15-second fallback
aborts the connection attempt if WebRTC fails.

\subsection{Call Initiation}
\label{sec:call-init}

\begin{enumerate}
    \item User clicks call or video-call button in chat header.
    \item \code{call\_request} emitted via Socket.io with \code{\{ callerId, calleeId, type \}}.
    \item Caller sees ``Calling...'' overlay; callee sees incoming call modal.
    \item Callee accepts: \code{call\_accepted} emitted; caller creates RTCPeerConnection.
    \item Caller gets media stream (audio only for voice, audio+video for video).
    \item Caller generates and emits \code{offer} (SDP).
    \item Callee sets remote description, generates \code{answer}, emits it.
    \item ICE candidates are exchanged via \code{ice\_candidate}.
    \item When \code{iceConnectionState === 'connected'}: call UI appears.
\end{enumerate}

\subsection{Call UI}
\label{sec:call-ui}

The call interface shows:
\begin{compactitem}
    \item Full-screen dark overlay.
    \item Remote video (or avatar + name for voice call).
    \item Local video as small PIP in bottom-right corner.
    \item Bottom bar: Mute, Camera toggle (video), End call.
    \item Timer at top showing call duration.
    \item Incoming call modal: ringing animation, accept/decline buttons.
\end{compactitem}

Call state is stored in-memory only (\code{CallContext}) and is lost on reload.

\subsection{Socket Events for Calls}
\label{sec:call-socket}

\begin{lstlisting}[style=js]
// Client -> Server
call_request, call_cancelled, call_accepted, call_rejected,
call_ended, offer, answer, ice_candidate

// Server -> Client (relayed to specific user)
incoming_call, call_cancelled, call_accepted, call_rejected,
call_ended, offer, answer, ice_candidate, user_busy
\end{lstlisting}

Active calls are tracked in an in-memory \code{Map}:
\code{activeCalls: callId -> \{ callerId, calleeId, type, startedAt \}}.

% --- SECTION 8: API REFERENCE ---
\section{API Reference}
\label{sec:api}

\subsection{Authentication Endpoints}
\label{sec:auth-endpoints}

\begin{table}[H]
    \centering
    \caption{Authentication API}
    \begin{tabular}{lllp{5.5cm}}
        \toprule
        \textbf{Method} & \textbf{Endpoint}         & \textbf{Auth} & \textbf{Description} \\
        \midrule
        POST & \code{/api/auth/register}   & None     & Create new account \\
        POST & \code{/api/auth/login}      & None     & Login (sets httpOnly cookies) \\
        POST & \code{/api/auth/logout}    & Required & Clear cookies \\
        POST & \code{/api/auth/refresh}   & None     & Refresh tokens \\
        \bottomrule
    \end{tabular}
\end{table}

\subsection{Message Endpoints}
\label{sec:message-endpoints}

\begin{table}[H]
    \centering
    \caption{Messaging API}
    \begin{tabular}{lllp{5.5cm}}
        \toprule
        \textbf{Method} & \textbf{Endpoint}          & \textbf{Auth} & \textbf{Description} \\
        \midrule
        GET  & \code{/api/messages/:userId} & Required & Get conversation history \\
        POST & \code{/api/messages/upload}   & Required & Upload image or file (multipart) \\
        \bottomrule
    \end{tabular}
\end{table}

\code{POST /api/messages/upload} returns:
\code{\{ url, filename, size, type \}}.

% --- SECTION 9: DATABASE SCHEMA ---
\section{Database Schema}
\label{sec:database}

\subsection{User Model}
\label{sec:user-schema}

\begin{lstlisting}[style=js]
const userSchema = new Schema({
    username:    { type: String, unique: true },
    email:       { type: String, unique: true },
    password:    { type: String, select: false },  // bcrypt
    avatar:      { type: String, default: '' },
    friends:     [{ type: Schema.Types.ObjectId, ref: 'User' }],
    friendRequests: [{ type: Schema.Types.ObjectId, ref: 'User' }],
}, { timestamps: true });
\end{lstlisting}

Password is excluded from queries by default (\code{select: false}).
Use \code{.select('+password')} when comparing.

\subsection{Message Model}
\label{sec:message-schema}

\begin{lstlisting}[style=js]
const messageSchema = new Schema({
    sender:      { type: Schema.Types.ObjectId, ref: 'User', required },
    recipient:   { type: Schema.Types.ObjectId, ref: 'User', required },
    content:     { type: String, default: '' },
    status:      { type: String, enum: ['sent', 'delivered', 'read'], default: 'sent' },
    roomId:      { type: String, index: true },  // conversation identifier
    replyTo:     { type: Schema.Types.ObjectId, ref: 'Message', default: null },
    attachments: [{
        type:     String,  // 'image' | 'file'
        url:      String,
        filename: String,
        size:     Number,
    }],
    reactions:   [{
        emoji:  String,
        userId: Schema.Types.ObjectId,
    }],
}, { timestamps: true });
\end{lstlisting}

\index{attachments}
\index{reactions}
\index{roomId}

Both \code{attachments} and \code{reactions} are recent additions to the schema
to support the media sharing module.

% --- SECTION 10: UI AESTHETIC ---
\section{UI Aesthetic}
\label{sec:ui}

\pingme{} uses a dark cyberpunk visual theme. Key design decisions:

\begin{compactitem}
    \item \textbf{Background}: Deep midnight navy (\code{#0D0B1E}).
    \item \textbf{Accents}: Electric violet (\code{#7B2FF7}), neon cyan (\code{#00D9FF}),
          neon pink (\code{#FF2D78}).
    \item \textbf{Fonts}: Space Grotesk (headlines), Inter (body),
          Material Symbols Outlined (icons).
    \item \textbf{Animations}: Framer Motion 3D flip transitions on AuthLayout
          between Login/Register. Custom CSS animations for message pop-in
          and fade-in effects.
\end{compactitem}

Tailwind v4 theme tokens are defined in \file{client/src/index.css} using the
CSS-first configuration style (no \file{tailwind.config.js} file exists).
Custom utilities include \code{prism-border}, \code{nebula-bg},
\code{animate-message-pop}, and \code{animate-fade-in}.

% --- SECTION 11: ROADMAP ---
\section{Development Roadmap}
\label{sec:roadmap}

\subsection{Current Status}
\label{sec:current-status}

\begin{table}[H]
    \centering
    \caption{Feature Completion Status}
    \begin{tabular}{llp{5cm}}
        \toprule
        \textbf{Feature}               & \textbf{Status} & \textbf{Notes} \\
        \midrule
        Authentication                & \color{success}{\textbf{Done}} & JWT, bcrypt, httpOnly cookies \\
        Real-time Messaging           & \color{success}{\textbf{Done}} & Socket.io relay, typing, read receipts \\
        Friend Management              & \color{success}{\textbf{Done}} & Request, accept, list, search \\
        UI/UX                         & \color{success}{\textbf{Done}} & Cyberpunk redesign, animations \\
        Media \& File Sharing         & \color{accent}{\textbf{In Progress}} & Image upload, lightbox, reactions, gallery \\
        Voice \& Video Calls          & \color{accent}{\textbf{In Progress}} & WebRTC, call UI, signaling \\
        Group Chat                    & \color{accentPink}{\textbf{Planned}} & Room management \\
        Push Notifications            & \color{accentPink}{\textbf{Planned}} & Service workers \\
        Message Encryption            & \color{accentPink}{\textbf{Planned}} & E2E encryption \\
        Stories / Status Updates      & \color{accentPink}{\textbf{Planned}} & Ephemeral content \\
        \bottomrule
    \end{tabular}
\end{table}

\subsection{Next Steps}
\label{sec:next-steps}

\begin{compactenum}
    \item \textbf{Complete media sharing}: finish lightbox, reaction badges,
          media gallery panel on the frontend.
    \item \textbf{Implement voice/video calls}: build \code{CallContext},
          \code{useWebRTC} hook, \code{IncomingCallModal}, and \code{CallOverlay}.
    \item \textbf{TURN server}: configure for production users behind symmetric NAT.
    \item \textbf{Group chat}: room-based architecture with Socket.io namespaces.
    \item \textbf{Push notifications}: service worker registration,
          FCM or Web Push integration.
    \item \textbf{Message encryption}: implement Signal Protocol or E2E encryption.
\end{compactenum}

% --- SECTION 12: KNOWN GAPS ---
\section{Known Gaps and TODOs}
\label{sec:gaps}

\begin{compactitem}
    \item ``Reject friend request'' button in \code{Sidebar.jsx} has no backend route.
    \item Persistent unread counts are client-side only (not stored in database).
    \item No test suite (unit, integration, or E2E).
    \item No backend ESLint or Prettier configuration.
    \item No Docker setup or CI/CD pipeline.
    \item \file{.env} files are tracked in git --- do not commit real secrets.
    \item Google OAuth is defined in requirements but not implemented.
    \item TURN server not configured (call quality affected behind symmetric NAT).
    \item Persistent call state is lost on page reload.
\end{compactitem}

% --- SECTION 13: CONCLUSION ---
\section{Conclusion}
\label{sec:conclusion}

\pingme{} demonstrates a complete full-stack real-time application architecture.
The server exposes a clean REST API alongside a Socket.io server for bidirectional
events, while the React frontend manages both REST and real-time state cleanly.

The most critical outstanding work is completing the media sharing module
(images, file cards, lightbox, reactions, gallery) and the WebRTC call
interface. Once these two features are complete, the application reaches
feature parity with modern chat applications at a fraction of the complexity
of commercial alternatives.

The codebase is well-structured, with clear separation between components,
contexts, hooks, and services. The cyberpunk UI theme provides a distinctive
visual identity that sets \pingme{} apart from generic messaging applications.

\vfill
\begin{center}
    \rule{0.7\linewidth}{0.5pt} \\
    \small\color{muted}
    \pingme{} Project Report --- April 2026 \\
    Built with React 19, Express 5, Socket.io 4, MongoDB, and Vite 7
\end{center}

\end{document}